\section{Discussão}

O problema é altamente paralelizável: cada ponto pode ser gerado e testado
sem depender dos demais. A comunicação ocorre apenas no final, quando os
acertos locais são somados. Essa baixa comunicação explica o speedup próximo
do ideal observado no teste.

Do ponto de vista de Flynn, a solução é MIMD, pois vários processos executam
instruções independentes sobre dados diferentes. O estilo de programação é
SPMD: todos executam o mesmo programa, mas cada \textit{rank} trabalha sobre
uma parcela dos pontos.

MPI adota o modelo de memória distribuída entre processos. Na máquina
multicore usada, a memória física é compartilhada, porém cada processo possui
seu próprio espaço de endereçamento. Essa abordagem também permite executar o
mesmo código em várias máquinas. Para este problema, MPI é adequado porque o
trabalho pode ser dividido com pouca troca de mensagens.

O speedup pode diminuir para poucos pontos, pois inicialização, sincronização
e redução passam a representar uma fração maior do tempo total. Também não se
espera que o ganho cresça indefinidamente devido à parte serial e aos recursos
limitados da máquina.
